% !TEX TS-program = xelatex
% File: house-style-and-preamble.tex
% Purpose: Brett's typical preamble + usage conventions:
%   examples (langsci-gb4e), italics vs. quotation, dashes, contractions, plain style, etc.

\documentclass[12pt]{article}

% --- Layout & language ---
\usepackage[margin=1in]{geometry}
\usepackage[british]{babel}            % British conventions; use -ize spellings in prose
\usepackage{fontspec}                  % Xe/LuaLaTeX
\setmainfont{Charis SIL}
\usepackage[final]{microtype}
\usepackage{marvosym}                  % \Cross symbol for cross-linguistic subscripts

% --- Quotation marks ---
\usepackage{csquotes}                  % \enquote{…} with locale-aware quoting
\usepackage{orcidlink}

% --- Hyperlinks ---
\usepackage{hyperref}
\hypersetup{
  colorlinks=true,
  linkcolor=blue,
  citecolor=blue,
  urlcolor=blue,
  pdfauthor={Brett Reynolds},
  pdftitle={House Style \& Preamble}
}

% --- Numbered linguistic examples (LangSci/gb4e wrapper, no 'exe' env) ---
\usepackage{langsci-gb4e}
\noautomath

% --- Lists & small utilities ---
\usepackage{enumitem}
\setlist{itemsep=0.3\baselineskip, topsep=0.3\baselineskip}
\usepackage{xspace}

% =========================
% Bibliography (biblatex)
% =========================
% Default portable setup:
\usepackage[backend=biber,style=apa,natbib=true,doi=true,isbn=false,url=true]{biblatex}
\addbibresource{references.bib}

% If working in LangSci projects, you can switch to their unified style:
% \usepackage[backend=biber,style=unified,natbib=true,doi=true,isbn=false,url=false]{biblatex}

% =========================
% Light house macros
% =========================
% Mention (italics) vs. quotation marks
\newcommand{\term}[1]{\textit{#1}}

% Small-caps abbreviations for glosses
\newcommand{\abbr}[1]{\textsc{#1}}

% Cross-linguistic subscript marker (e.g., \textsc{subject}\crossmark)
\newcommand{\crossmark}{\textsubscript{\Cross}}

% Judgement markers
\newcommand{\ungram}[1]{*\!#1}
\newcommand{\marg}[1]{?\!#1}
\newcommand{\odd}[1]{\#\!#1}

% e.g., i.e., etc., with sensible spacing
\newcommand{\eg}{e.g.,\xspace}
\newcommand{\ie}{i.e.,\xspace}
\newcommand{\etc}{etc.\xspace}

% =========================
% Document
% =========================
\title{House Style \& Preamble:\\A Compact Demonstration}
\author{Brett Reynolds \orcidlink{0000-0003-0073-7195}\thanks{I used ChatGPT 5, Claude Sonnet 4.5, and Gemini Pro 2.5 extensively in drafting and revising the paper. I reviewed, edited, and approved all the material and take full responsibility for the final text and conclusions. \href{mailto:brett.reynolds@humber.ca}{brett.reynolds@humber.ca}}\\
Humber Polytechnic \& University of Toronto}
\date{\today}

\begin{document}
\maketitle

\section{Mention vs.\ quotation} \label{mention}
Use italics for mention of forms and categories; reserve quotation marks for actual quotations.
Mention: \term{going forward} is best analysed as a preposition selecting a directional PP.
Quotation (with \verb|\enquote|): As \enquote{going forward} spread in business English, its syntax broadened.
Nested quotes: \enquote{Outer with an \enquote{inner} quote}.
Semantic meanings: use single `quote'

\section{Dashes, ranges, hyphens} \label{dashes}
Use an en dash (~-- ) with surrounding spaces for parenthetical asides: the category is stable~-- within limits~-- across registers.
Use an en dash for ranges and relations: 2001--2025; pp.\ 113--127; the form--function distinction.
Use hyphens only for compounds: corpus-based study, task-specific rubric.

\section{Contractions and plain style} \label{contractions}
Contractions are preferred: \enquote{We don’t gain anything by multiplying categories.}
Prefer direct verbs and short clauses. Delete throat-clearers:
\begin{itemize}[nosep]
  \item \textbf{Avoid:} \enquote{It is important to note that the results clearly show that \dots}
  \item \textbf{Prefer:} \enquote{The results show \dots}
  \item \textbf{Avoid:} Overly strong modals like \enquote{must}. E.g., \enquote{The features must cluster reliably.}
  \item \textbf{Prefer:} Softer phrasing. E.g., \enquote{The features should cluster reliably.}
\end{itemize}
Prefer simple coordinators (and, but) to hackneyed "academic" adverbs (moreover, furthermore, nevertheless, yet)
Sentence-initial coordinators are fine where they improve flow. Avoid overuse.

\section{Numbered examples with \texttt{langsci-gb4e}}\label{numbered-examples}
Single example (note: no \texttt{exe} environment):
\ea\label{ex:simple}
\textit{The committee have decided to adjourn.}
\z

Judgement markers and subexamples:
\ea\label{ex:sub}
	\ea \textit{She has already left.}
	\ex[*]{\ungram{\textit{She have already left.}}}
	\ex[\#]{\odd{\textit{The square triangle laughed.}}}
	\z
\z

Cross-reference: see (\ref{ex:sub}).

Interlinear gloss (\texttt{\textbackslash gll}/\texttt{\textbackslash glt}; abbreviations in small caps):
\ea\label{ex:gloss}
\gll Ich sehe den Hund. \\
     I see  the.\abbr{acc} dog \\
\glt `I see the dog.'
\z

\section{Category vs.\ function} \label{category}
Distinguish lexical categories (e.g., \term{determinative}, \term{adjective}) from functions (e.g., \term{head}, \term{modifier}, \term{subject}).
In \enquote{\term{those} books}, \term{those} is a \term{determinative} (category) functioning as a \term{determinative} (function) of the NP.
Use the cross symbol to mark comparative concepts: \textsc{subject}\crossmark, \textsc{topic}\crossmark, \textsc{noun}\crossmark. Avoid literal \texttt{\textsubscript\{cross\}} in new writing.

\section{Minimal citation examples (biblatex)}\label{citation}
Parenthetical: \citep{HuddlestonPullum2002}. Textual: \textcite{Pullum2018}.
Multiple: \citep{HuddlestonPullum2002,Pullum2018}. Page: \citep[113--127]{HuddlestonPullum2002}.

\section{Journal article formatting}\label{journal-formatting}
Scholarly journals use structural apparatus sparingly. Avoid textbook-style formatting.

\subsection{Sections and subsections}
\textbf{Use:} \verb|\section{}| and \verb|\subsection{}| for major divisions. Keep subsections to substantial content (3+ paragraphs or equivalent). A one-paragraph subsection should be integrated into the section prose.

\textbf{Avoid:} \verb|\paragraph{}| headings. Replace with strong topic sentences and discourse markers.

\begin{itemize}[nosep]
  \item \textbf{Wrong:} \verb|\paragraph{Flexible lexical categories.}| What about languages like Salish...
  \item \textbf{Right:} A first objection concerns languages with flexible lexical categories. What about Salish...
\end{itemize}

\subsection{Bold, bullets, enumeration}
\textbf{Avoid bold labels in prose.} Taxonomies and argument sequences should flow as narrative, not as itemized lists.

\begin{itemize}[nosep]
  \item \textbf{Wrong:} \textbf{(i) Cross-linguistic property covariation.} The diagnostic features must cluster...
  \item \textbf{Right:} First, the diagnostic features must cluster...
\end{itemize}

\textbf{Enumerated lists} are acceptable for:
\begin{itemize}[nosep]
  \item Hypothesis or prediction lists (Section 9 in main.tex)
  \item Linguistic examples (\verb|\ex| from langsci-gb4e)
  \item Illustrative codebook entries
\end{itemize}

\textbf{Avoid} for:
\begin{itemize}[nosep]
  \item Argumentative sequences (use prose with transitions: ``first,'' ``second,'' ``finally'')
  \item Objection-reply pairs (use discourse markers: ``A first objection,'' ``A related worry'')
  \item Before/after pedagogical contrasts (use narrative: ``Traditional accounts treat... The present framework dissolves...'')
\end{itemize}

\subsection{Transitions and discourse markers}
Use ordinal markers and discourse connectives for flow:
\begin{itemize}[nosep]
  \item \textbf{Sequence:} first, second, third; one [concern], another, a final [point]
  \item \textbf{Objections:} A first objection concerns...; A second question asks...; A related worry...; Finally, does...
  \item \textbf{Illustrations:} Consider first...; A second illustration...; The clearest case...
\end{itemize}

\section{Dual-audience accessibility}\label{dual-audience}
When writing for readers from multiple disciplines (e.g., linguists and philosophers), make cross-disciplinary terminology accessible without sacrificing rigor.

\subsection{Paragraph length}
Keep paragraphs under 100 words on average. Dense 200+ word paragraphs should be broken into 2--4 shorter units with clear topic sentences. Aim for ~60 words per paragraph in body text.

\subsection{Parenthetical glosses}
When introducing technical terms, add brief functional explanations:
\begin{itemize}[nosep]
  \item Latent variables (quantities representing unobserved category strength)
  \item Spearman's $\rho$ (a rank correlation measure)
  \item Differential object marking (where only some objects trigger case or agreement)
\end{itemize}

\subsection{Concrete before abstract}
Introduce technical apparatus with concrete examples or plain-language motivation:
\begin{itemize}[nosep]
  \item \textbf{Wrong:} Define a vector of observable diagnostics: $\mathbf{n}_L = \langle \text{ArgHead}, \dots \rangle$.
  \item \textbf{Right:} Traditional typology trades in binary tallies: language X ``has'' adjectives or it doesn't. The comparanda~$\times$~categories matrix provides something better: explicit measurement models. For nominality, define a vector...
\end{itemize}

\subsection{Comparanda notation and terminology}
Maintain a strict distinction between comparative concepts and language-specific categories:
\begin{itemize}[nosep]
  \item Use \textsc{Term}\textsubscript{cross} for cross-linguistic comparanda (syntactic functions, semantic targets, discourse roles, comparative categories). Note: always use \verb|\textsubscript{}| for these labels; do not use the underscore character (\_) directly in text mode.
  \item Use \textsc{Term}\textsubscript{\text{eng}} (or other language tag) for language-specific realisations; \textsc{Term}\textsubscript{\text{int}} for language-internal discussion without naming a language.
  \item Reserve \emph{function} (without modifier) for syntactic functions only; use \emph{target} for semantic targets and \emph{role} for discourse/pragmatic roles. Use \emph{category} for lexical lexical categories or their phrasal projections.
  \item Remind readers that mappings between axes are rarely one-to-one: a single expression can realise multiple functions, targets, and roles simultaneously.
\end{itemize}

\section{BibTeX conventions}\label{bibtex}
Protect capitals in proper nouns and first words after colons using braces:

\begin{verbatim}
title = {The {Cambridge} Grammar of the {English} Language}
title = {Radical Construction Grammar: {Syntactic} Theory...}
title = {Character identity mechanisms: {A} conceptual model...}
\end{verbatim}

\textbf{What to protect:}
\begin{itemize}[nosep]
  \item Proper nouns: \{Cambridge\}, \{English\}, \{American\}
  \item First word after colon: \{Syntactic\}, \{A\}, \{An\}
  \item Acronyms: \{HPC\}, \{DOM\}
\end{itemize}

\textbf{What NOT to protect:}
\begin{itemize}[nosep]
  \item Common nouns, adjectives, verbs (unless part of proper noun)
  \item Title-initial words (BibTeX handles these automatically)
\end{itemize}

\clearpage
\printbibliography

\end{document}
